\lecture{6}{Wed 14 Sep 2022 10:32}{}

\begin{lemma}
	If $\varphi:G \to  G'$ is a homomorphism of groups, then $\varphi(G)\le G'$.
\end{lemma}
\begin{proof}
	Homework exercise.
\end{proof}

\begin{theorem}
	If $\varphi:G \to G'$ is a homomorphism of groups, then: 
	\begin{enumerate}
		\item $\operatorname{ker} \varphi \le G$
		\item For all $g \in G$, one has $g^{-1} \left( \operatorname{ker} \varphi \right) g \subset \operatorname{ker} \varphi$
	\end{enumerate}
\end{theorem}
\begin{proof}
	Suppose $a,b \in \operatorname{ker} \varphi$. Then \[
		\varphi(ab^{-1} )=\varphi(a)\varphi(b)^{-1} = e' (e')^{-1} = e'
	.\] 
	So $ab^{-1}  \in \operatorname{ker} \varphi$. Thus $\operatorname{ker} \varphi \le G$ by subgroup criterion.
	Whenever $k \in \operatorname{ker} \varphi$, one has for all $g \in G$, \[
		\varphi(g^{-1} kg) = \varphi(g)^{-1} \varphi(k)\varphi(g) = \varphi(g)^{-1} e' \varphi(g) = \varphi(g)^{-1} \varphi(g) = e'
	.\] 
\end{proof}

\begin{remark}
	The kernel is \textbf{fixed under conjugation}, and forms a \textit{normal subgroup}. Also note: we actually have $g^{-1} \left( \operatorname{ker} \varphi \right) g = \operatorname{ker} \varphi$.
\end{remark}

\begin{definition}[Normal Subgroup]
	Suppose $N\le G$, then $N$ is a \textit{normal subgroup} of $G$ when \[
		g^{-1} N g \subset N
	\]
	for all $g \in G$. We then write $N \triangleleft  G$.
\end{definition}

If $g^{-1} Ng \subset N$, then $g^{-1} Ng = N$ (exercise). Then $N = g N g^{-1}$. So if these relations hold for all $g \in G$, we have \[
	g N = N g
\] for all $g \in G$.

\begin{example}
	\begin{enumerate}
		\item If $G$ is abelian and $H\le G$ then $H \triangleleft G$. Observe that for all $G \in G$,and $h \in H$, one has $g^{-1} hg=hg^{-1} g = he = h$, so $g^{-1} Hg = H$ for all $g \in G$.
		\item Recall $Z(G) = \{z \in G: \forall x \in G, z x = xz\} $. Thus, for all $g \in G$, one has $g^{-1} Z(G) g = Z(G)$. Hence $Z(G) \triangleleft G$.
		\item Consider the group
			\[
				G = D_{2n} = \left\langle \sigma, \tau: \sigma^n = \tau^2, \sigma\tau = \tau\sigma \right\rangle
			.\]
		Consider $H = \left\langle \sigma \right\rangle $. We check that $H\triangleleft G$. We have \[
			\sigma^{-m} (\sigma^{l})\sigma^{m}= \sigma^{l} \in H.
		.\] 
		\begin{align*}
			(\tau\sigma^m)^{-1}\sigma^l(\tau\sigma^m) &= \sigma^{-m}\tau^{-1} \left( \sigma^l \tau \right)  \sigma^{m}\\
			 &= \sigma^{-m}\tau^{-1} \left( \tau \sigma^l \right)  \sigma^{m}\\
			 &= \sigma^{-m}\sigma^{-l}\sigma^m \\
			 &= \sigma^{-l} \in H 
		.\end{align*} 
		Thus $g^{-1} H g \subset H$ for all $g \in D_{2n}$. Hence $H \triangleleft D_{2n}$.
	\end{enumerate}
\end{example} 


\subsection{Quotient Groups (or Quotient Groups)}

\[
	(\mathbb{Z}_n, +) \cong \mathbb{Z} / n \mathbb{Z}
.\] 

Recall that when $H \le G$, we can define $a\sim b$ when $ab^{-1} \in H$. Then $[a] = Ha$. 

\begin{theorem}
	Suppose $N \triangleleft G$. Define \[
		G / N = \{Na: a \in G\} 
	.\] 
	Then $G / N$ is a group relative to the group operation \[
		Na \cdot Nb = N(ab)
	.\] 
\end{theorem}
\begin{proof}
	\begin{itemize}
		\item (well-defined group operation) If $Na = Na'$ and $Nb = Nb'$, then $Na \cdot Nb = Na' \cdot Nb'$.

		If $Na = Na'$, then $a' = n_1a$ for some $n_1\in N$. And if  $Nb=Nb'$, then $b'=n_2b$ for some $n_2\in N$.
		Thus
		\begin{align*}
			a'b'&=n_1an_2b\\
			    &=n_1(an_2a^{-1} )ab \\
			    &\in n_1(aNa^{-1} )ab\\
			    &= n_1Nab \\
			    &= Nab \\
		\end{align*}
		Hence $Na'b'  \subset Nab$. By symmetry $Nab \subset Na'b'$, hence $Nab = Na'b'$. 
		Thus the putative group operation on $ G / N$ is well-defined.

		\item (closure)
		follows from definition of the product.
		\item (associativity)
		inherited from the group $G$. \[
			(Na)(Nb Nc) = Na(Nbc) = Na(bc) = N(ab)c = (Nab)Nc = (Na Nb) Nc
		.\] 
		\item (identity)
			$Ne = N$ serves as identity:
			\[
			Na Ne = Nae = Na = Nea = Ne Na
			.\] 
		\item (inverse)
			$Na^{-1} $ serves as inverse of $Na$. Similar argument.
	\end{itemize}
	So $G / N$ is indeed a group.
\end{proof}

